\section{Additional Figures}\label{sec:additional_figures}
In \cref{sec:thorough_analysis} we enumerate all undetected fault configurations
in the distance-3 and -5 double checks.
\Cref{fig:fcc_against_fault_count} shows
the number of such configurations,
after grouping by fault count and whether they are benign or malignant.
\begin{figure}[t]
	\centering
	\includegraphics[width=0.72\linewidth]{figures/fcc_against_fault_count.pdf}
	\caption{The number of undetected fault configurations
	in the double-check circuit for magic state cultivation,
	grouped by fault count i.e.\ number of faults in the configuration.
	(a) and (b) correspond to distance-3 and -5 cultivation respectively.
	In both subfigures:
	blue (orange) data correspond to cultivating S (T) states.
	Benign (malignant) configurations
	lead to a logical fidelity of 1 (0).
	Five low-fault-count malignant counts are annotated.}
	\label{fig:fcc_against_fault_count}
\end{figure}

In \cref{sec:z_flagged_double_checks} we show
our Z-flagged distance-3 double check as a detector slice diagram.
\Cref{fig:d3a6f2_conventional_diagram} shows
the same circuit more conventionally.
\begin{SCfigure}
	\begin{tikzpicture}
\begin{yquant}

qubit {0} q0;
qubit {a} qa;
qubit {b} qb;
qubit {1} q1;
[darkred]
qubit {\alpha} qalpha;
qubit {c} qc;
qubit {2} q2;
qubit {d} qd;
qubit {3} q3;
[darkred]
qubit {\beta} qbeta;
qubit {4} q4;
qubit {e} qe;
qubit {5} q5;
qubit {6} q6;
qubit {f} qf;

setstyle {darkred, dashed} qalpha,qbeta;
discard qa,qb,qalpha,qbeta,qc,qd,qe,qf;
hspace {1mm} -;
box {T} q1,q4,q5;
box {T$^\dag$} q0,q2,q3,q6;
init {$\ket{+}$} qa,qb,qc,qd,qe,qf;
align -;

cnot q0|qa;
cnot q4|qe;
cnot q3|qd;
cnot q2|qc;
cnot q1|qb;
cnot q6|qf;
[darkred]
init {$\ket{0}$} qalpha,qbeta;

[dashedcnot]
cnot qalpha|q1;
[dashedcnot]
cnot qbeta|q2;
[dashedcnot]
cnot qalpha|q2;
[dashedcnot]
cnot qbeta|qe;

align -;
cnot qa|q1;
cnot qd|q2;
cnot qf|qe;

cnot q1|q2;
cnot q5|qe;

cnot qe|q2;

dmeter {X} q2;
discard q2;
init {$\ket{+}$} q2;
align q2;

cnot qe|q2;

align -;
cnot q1|q2;
cnot q5|qe;

cnot qa|q1;
cnot qd|q2;
cnot qf|qe;

[dashedcnot]
cnot qalpha|q2;
[dashedcnot]
cnot qbeta|qe;
[dashedcnot]
cnot qalpha|q1;
[dashedcnot]
cnot qbeta|q2;

align -;
cnot q0|qa;
cnot q4|qe;
cnot q3|qd;
cnot q2|qc;
cnot q1|qb;
cnot q6|qf;
[darkred]
dmeter {Z} qalpha,qbeta;
discard qalpha,qbeta;

align -;
dmeter {X} qa,qb,qc,qd,qe,qf;
discard qa,qb,qc,qd,qe,qf;
box {T$^\dag$} q1,q4,q5;
box {T} q0,q2,q3,q6;

hspace {1mm} -;

\end{yquant}
\end{tikzpicture}

	\caption{A Z-flagged version of the distance-3 double-check circuit
	that detects the four fault configurations
	shown in \cref{fig:all_2_fault_configurations}.
	The two added flag qubits {\alpha} and {\beta},
	and the four additional CX layers,
	are drawn red and dashed.}
	\label{fig:d3a6f2_conventional_diagram}
\end{SCfigure}